\section{Inference Systems Deep Dive}
\sectionkicker{MEMORY / PREFETCH / DECODING}
\begin{wideflow}
{\Large\bfseries 推論を速くする場所は、メモリか、先読みか、復号か}\par
\smallskip
{\color{SurveyMuted}vToken、OasisKV、Ripple-Pivot Searchは、同じ「高速化」でも別のボトルネックを動かす。KVの置き場所と生存管理、prefetch、diffusion LLMのdecoding policyを共通の問いで比較する。}\par
\end{wideflow}

三つの提案は、推論の別々の詰まりを狙う。

\begin{themeoverview}[機構を同じ表に置く]
\begin{tabularx}{\linewidth}{@{}>{\bfseries}p{0.27\linewidth}X@{}}
Memory placement / liveness & vTokenはlogical tokenの生存とphysical KV blockをindirectionで分離し、非同期repackingを使う。paperはPagedAttentionとCUDA Graphとの互換性も記述する\cite{vtoken}。\\
Tiering / prefetch & OasisKVはHBMにはrelevantなKVだけを置き、より容量の大きいmemory tierを使いながらlookahead sparse prefetchingを行う\cite{oasiskv}。\\
Decoding policy & Ripple-Pivot Searchはmid-entropy pivotとlookahead evaluationを使うtraining-freeのdiffusion-LLM decoding methodである\cite{ripplepivot}。
\end{tabularx}
\end{themeoverview}

比較のポイントは、三つを一つの性能ランキングにすることではない。vTokenとOasisKVはKVをどう生かし、どこに置き、どう先読みするかを変える。一方Ripple-Pivotは、diffusion LLMが次のtokenを決めるdecoding policyそのものを変える。前者はmemory systemの配置・再利用、後者は生成手順の探索を動かすため、「速くする」の中身が異なる。

\begin{claimboundary}[評価値の帰属]
OasisKVはvLLM上のaccuracyとthroughput、Ripple-Pivotは三つのdLLMと四つのbenchmarkでspeedとqualityの結果を報告するが、今回確認できる範囲では、いずれもpaper authorsが報告した測定である。独立再現結果や横断的な同一条件の比較として言い換えない。
\end{claimboundary}

読者が「推論を速くする」と聞いたら、memory placement、prefetch、decoding policyのどこを変えたのかを先に確認したい。そうすれば、KV cacheの運用を組み替える提案と、復号の探索を変える提案を、同じ指標だけで混同せずに済む。次の観測点は、各機構がどのruntime・モデル・評価条件で有効なのかを、論文の報告境界ごとに追うことである。
